\section{Introduction}
\label{sec:intro}

An estimated 536.6 million adults aged 20--79 years were living with diabetes in 2021, and this number is projected to reach 783.2 million by 2045~\cite{sun2022idf}. Type~2 diabetes mellitus (T2DM) accounts for the large majority of cases, and a substantial share of people with diabetes are undiagnosed~\cite{idf2021}. Because hyperglycaemia damages the kidneys, retina, nerves and vasculature years before symptoms appear, risk estimates from routinely collected measurements could help prioritise confirmatory testing and follow-up.

Machine-learning (ML) classifiers trained on variables such as plasma glucose, body-mass index (BMI), blood pressure and familial risk are a common route to such risk-prediction tools. Two design decisions dominate this literature. The first is feature selection: wrapper methods such as recursive feature elimination (RFE)~\cite{guyon2002rfe,kohavi1997wrappers} prune redundant inputs, and RFE with cross-validation (RFECV) chooses the subset size by validated performance. The second is model choice: stacked generalisation~\cite{wolpert1992stacked} trains a meta-learner on the out-of-fold predictions of diverse base learners and is often reported to outperform any single learner.

On small clinical datasets, however, the reported gains of such pipelines are hard to interpret. Imputation, scaling, oversampling or feature selection are frequently fitted on the complete dataset before cross-validation or splitting. This leaks information from the evaluation data into model development and inflates performance~\cite{ambroise2002selection,kaufman2012leakage,kapoor2023leakage}. Moreover, differences between models are usually reported without confidence intervals or paired tests, calibration is rarely assessed~\cite{vancalster2019calibration}, and the default 0.5 decision threshold is seldom justified for the intended use. Systematic comparisons have also found no average performance benefit of flexible ML over logistic regression for clinical prediction~\cite{christodoulou2019}.

This study therefore asks a deliberately narrow question: \emph{when every data-dependent step is kept inside the validation loop, does combining RFECV with a heterogeneous stacking ensemble improve T2DM prediction on the PIMA cohort over simpler models, and how do calibration and the choice of operating threshold affect its interpretation?} The contribution is a rigorous evaluation framework, not a new learning algorithm:

\begin{enumerate}
  \item A leakage-controlled pipeline in which median imputation of physiologically impossible zeros, clinically motivated feature engineering, standardisation, RFECV and an eight-learner stacking ensemble are all re-fitted within the training part of every split.
  \item A common feature-selection protocol in which RFECV is fitted once per outer fold and its subset is shared by every compared learner, together with a learner-specific RFECV sensitivity analysis for logistic regression.
  \item A statistically explicit evaluation: 5$\times$10-fold repeated cross-validation with corrected confidence intervals and corrected resampled $t$-tests~\cite{nadeau2003inference,bouckaert2004replicability}, Holm adjustment for multiplicity~\cite{holm1979}, bootstrap confidence intervals and DeLong tests~\cite{delong1988,sun2014delong} on an untouched hold-out set.
  \item A clinically oriented assessment that goes beyond discrimination: calibration, pre-specified sensitivity-targeted thresholds, decision-curve analysis~\cite{vickers2006dca}, a missing-data sensitivity analysis and SHAP explanations~\cite{lundberg2017shap}.
\end{enumerate}

We report where the framework helps and where it does not. In particular, under strict validation and prespecified hyperparameters, the stacking ensemble does not discriminate better than regularised logistic regression on this cohort. Because PIMA predictors include 2-h oral glucose tolerance test (OGTT) values, the setting is risk stratification among people who have already had an OGTT, not population-level screening, and the analysis is primarily a methodological benchmark.
